\title{Complex fourfolds from a marked rank-six torus family}
\author{AI-produced mathematical text}
\date{18 November 2026}
\maketitle

\begin{abstract}
The marked rank-six monodromy with finite orders five and four and its
primitive maximal cusp produces compact complex fourfolds on
$3\mathbb{HP}^2\#(-\mathbb{HP}^2)$. Their explicitly normalized family
is parametrized by the half-plane $\operatorname{Im}\xi>c_*$, where
$c_*=1.0945402503114027801162682630549457037\ldots$ is defined by
a certified period minimum. Its members are pairwise nonisomorphic,
have algebraic dimension one and automorphism group
$\mathbb C\times\mathbb Z/2$, and have no antiholomorphic
isomorphisms between them. Wall's extended-form classification and
Kramer--Stolz's 2007 homotopy-sphere absorption theorem supply the
smooth classification input. The calculations here give the actual
integral cusp maps, the genus-two exterior-square period construction,
the family invariants and its endpoint obstruction, together with an
effective classification of the retained integral inputs. The hard
topological step is the coefficient-one torsion boundary map; the
analytic step is the exact real-transversality minimum. Historical
checks are recorded by statement. GPT-6 Astra Pro re-derived every
retained proof in this writing and reran the stated exact certificates;
no human proof reading is recorded in this writing.
\end{abstract}

\paragraph{Known inputs and the precise additional calculations.}
The characteristic index polynomials are used in the explicit
almost-complex form explained in \S\ref{sec:V-known-inputs}.
The final smooth comparison uses the extended Wall form
\cite[Definition 3.6 and Theorem 3.7(1), pp.~12--14]{IVVWallForm}
and the absorption theorem
\cite[Theorem A, p.~178]{IVVKramerStolz}.
The open-base period construction uses the section form of the Oka
property \cite[Definition 7.7, p.~22; Proposition 7.13(d), p.~24]{IVVOkaOneV5}
after constructing the required local sprays. Each application retains
its marked lattice, source plane, finite quotient and cusp maps.

\paragraph{Reading the verification record.}
The current mathematical review covers all retained statement proofs,
including the complete cellular reduction in \cref{night:P1.19}.
Each proof's record distinguishes this review from earlier checks.
The current source openings identify the exact version and locator
used, and the computation records distinguish executed programs from
historical certificates. Typesetting and visual inspection have their
own records and confer no additional mathematical check.
